Pour chaque entier strictement positif $n$, on note $f(n)$ le nombre de façons d'écrire $n$ comme une somme de puissances de $2$ à exposants entiers positifs ou nuls. Deux représentations qui ne diffèrent que par l'ordre de leurs termes sont considérées comme identiques. Par exemple, $f(4) = 4$, car $4$ peut s'écrire des quatre façons suivantes : $4$ ; $2+2$ ; $2+1+1$ ; $1+1+1+1$.

Montrer que, pour tout entier $n \ge 3$, on a
\[ 2^{n^2/4} < f(2^n) < 2^{n^2/2}. \]
